\documentclass[11pt]{article}
\usepackage[margin=0.85in]{geometry}
\usepackage{amsmath,booktabs,xurl,hyperref}
\usepackage[T1]{fontenc}
\setlength{\emergencystretch}{2em}
\title{Why the Flow Seed Produces a Long Optimized Completion Tail}
\author{Pan Dynamics Collision Avoidance Research}
\date{October 1, 2026}
\begin{document}
\maketitle

\section{Finding}
The selected flow seed is already a fully admitted avoidance trajectory. The
long tail is primarily caused by its late admission under the current sufficient
all-future geometric certificate, then retained as a large optimization horizon.
The seed is useful for safety feasibility but is not simultaneously a good
nominal-recovery trajectory. Multiple solver calls also arise from the current
CLF/correction sequence; their count is not a count of failed flow initializations.

A diagnostic dispatcher exposes the current, unmodified local functions of
\path{controller/solvePredictiveControl.m} outside the production tree. At the
original 8-m/s \texttt{turningCrossing} initial state, its selected side-$-1$
flow has 216 holds and initial CLF slack 1.1994339752192968, matching the saved
original initialization. It takes 0.153186 seconds to generate in this probe,
excluding construction of the cached terminal core. Complete evaluation gives
zero hard residual, zero safety slack and all nominal interpolation intervals
certified. The minimum evaluated prediction collision margin is 0.757297 m.
This is a prediction result, not a continuous-plant certificate.

\section{Which terminal requirement keeps the seed growing?}
The minimum requested length is 68 holds, or 3.4 seconds: eight prefix holds
plus the existing three-second completion allowance. The seed grows until both
its intrinsic terminal-core membership and its all-future separation test pass.
The measured gates along this one selected seed are:
\begin{center}
\begin{tabular}{rrrrr}
\toprule
Time (s) & Hold & Core residual & Geometry margin (m) & Current gap (m) \\
\midrule
3.40 & 68 & 12.717761 & -30.154361 & 18.972105 \\
4.65 & 93 & $-3.6494\,10^{-9}$ & -29.547834 & 24.971709 \\
10.75 & 215 & $-1.4901\,10^{-8}$ & -0.373690 & 44.518837 \\
10.80 & 216 & $-1.4901\,10^{-8}$ & 0.020179 & 45.443667 \\
\bottomrule
\end{tabular}
\end{center}
Negative core residual means membership is satisfied. Its first admission is
at 4.65 seconds, while all-future geometry first passes at 10.80 seconds.
Thus, along this seed, an additional 6.15 seconds after core entry is caused
by the geometric certificate. The core itself is small: its dimensionless
weighted radius is 0.0001220703125, not a physical distance.

The target has constant sideslip 0.05 rad and acceleration 1 m/s$^2$, so its
center follows a fixed circle even as its angular speed changes. The circle is
centered at approximately $(44.813337,0)$ m, with radius 32.013337 m. The current
straight-terminal branch encloses the complete target body orbit in a disk of
radius 33.165428 m and requires the entire future ego ray, with body/core
padding, to avoid that disk. It discards simultaneous timing and also fills
the orbit's interior. A large actual same-time gap therefore does not make this
sufficient test pass. Failure of this lower bound is not proof of a collision.

The seed's first accepted endpoint fixes its array length. Later calls shift
that array by one hold, so 216 becomes 212 at frame 5. The sequential programs
use \texttt{size(anchor,2)}; they do not minimize horizon length or search all
earlier truncations of an improved plan. The 10.8-second seed endpoint is not
a proof that every feasible maneuver needs 10.8 seconds. Optimizing every
input in that continuation on every frame is an implementation choice, not
by itself a requirement of recursive-feasibility theory.

\section{What the flow initialization does and does not provide}
The initializer follows temporary moving-flow guidance, then switches its
internal rollout to straight terminal-trim feedback. The terminal pose is free;
its policy need not point along the given path. For the selected seed, heading
settles near $-24.7546^\circ$. Its 10.8-second endpoint has lateral position
$-20.386842$ m. This is the \emph{initial planned tail}, not the executed
closed-loop vehicle trajectory. The completed campaign subsequently optimizes
and recovers nominal behavior.

Consequently, the seed is good enough to bypass safety restoration but still
leaves substantial work for the CLF and nominal tracking objectives. The
initialization problem, the terminal certificate and the dissipation objective
are not computationally well aligned. A diagnostic opposite-bias seed ends at
204 holds but has a sampled future collision and hard residual 2.308035;
it is not an admissible shorter substitute. The original algorithm stops seed
screening after the first admitted seed, so this second seed was evaluated only
for the diagnosis. No new online maneuver mode was added.

\section{Why several optimizations remain}
Original frame 5 uses the shifted previous solution, not a newly generated
flow seed. Its seven numerical calls comprise four feasibility-correction QPs
and three CLF-related SOCPs, as established in
\path{CLF_SLOW_FRAME_20261001.tex}. The three conic calls try zero CLF relaxation,
then minimize positive relaxation after that local slice is infeasible, then
refine predictive cost while fixing the minimizing first input. Safety slack
is already zero, but final CLF slack remains 4.555671. Continuing CLF work is
consistent with the intended hierarchy.

The correction machinery exists because changing the first input changes the
entire nonlinear suffix. A shifted plan must also be reevaluated when the
measured state differs from its prediction. In the independent ODE45 experiment,
frame 5 differs from the preceding RK4 successor by only about $1.6\,10^{-9}$ m
in position and $3.25\,10^{-7}$ rad/s in yaw rate. This is integrator discrepancy,
not simulated observation error. The implementation's exact-equality check
therefore does not reuse the exact shifted cache. Long open-loop propagation
and a small terminal core can amplify sensitivity, and feasibility repair may
be needed after CLF changes. Original per-QP cause/timing records were not saved,
so this diagnosis does not assign all four QPs solely to that discrepancy.

The current organization repeatedly improves the first input, repairs later
controls, and resolves a large conic problem. A better seed can reduce that
work, but cannot by itself remove the conservative endpoint geometry or the
fixed inherited horizon. The design issues are the strong geometric closure,
the seed's off-path terminal continuation, and repeated processing of all tail
inputs. Future simplification must preserve a valid time-aligned continuation
certificate; simply deleting the tail after the first encounter is not justified
for this recurring circular-target fixture. Production algorithms, constraints
and stopping rules were not changed during this diagnosis.
\end{document}
